\section{风险评估难题：为何现有 Agent Benchmark 不够}

讲者指出，当前 Agent 评测普遍存在四个问题：缺乏标准化、开放性不足、可复现性差、集成成本高。很多 benchmark 沿用模型中心评测思路，内置固定 harness，导致新 Agent 接入成本高，跨 benchmark 对比困难。

\begin{importantbox}{N * M 集成困境}
若有 $N$ 个待测 Agent、$M$ 个 Benchmark，传统模式下往往需要近似 $N \times M$ 次集成适配。这不仅成本高，而且容易引入实现偏差，导致评测结果难以横向比较。
\end{importantbox}

\begin{table}[H]
\centering
\begin{tabular}{p{2.6cm}p{3.8cm}p{4.6cm}}
\toprule
\textbf{现状问题} & \textbf{直接后果} & \textbf{对研究/工程的影响} \\
\midrule
接口不统一 & 每个 benchmark 单独适配 & 评测成本高，迭代慢 \\
封闭 harness & 难替换组件与环境 & 创新难复现，结论难验证 \\
指标不一致 & 同名指标语义不同 & 排行可比性弱 \\
缺少对抗评测 & 只测正常任务性能 & 安全能力被高估 \\
\bottomrule
\end{tabular}
\caption{当前 Agent 评测体系的结构性问题}
\end{table}

\begin{knowledgebox}{风险评估应覆盖的最小维度}
\begin{itemize}
    \item 正常任务性能（utility）；
    \item 对抗鲁棒性（robustness）；
    \item 安全后果强度（impact）；
    \item 触发成本与攻击可扩展性（attack economics）；
    \item 复现性与可审计性（reproducibility/auditability）。
\end{itemize}
\end{knowledgebox}

\subsection{本章小结}
没有标准化评测，就没有可信比较；没有可信比较，就无法指导防御投入。Agent 安全首先需要评测基础设施。

